\lecture{9.7}{Front-Door Identification}
\label{lec:9-7}

\subsection{Front-Door Graph}

The front-door model shows that an effect can be identified even when there is
unmeasured confounding between treatment and outcome and no valid adjustment
set. The key is a measured mediator $M$ through which all effects of $A$ on $Y$
pass.

\begin{figure}[H]
\centering
\begin{tikzpicture}[
  node distance=1.35cm,
  box/.style={draw, rounded corners=2pt, minimum width=.9cm, minimum height=.7cm},
  arr/.style={-{Latex[length=2.3mm]}, thick}
]
\node[box] (A) {$A$};
\node[box, right=of A] (M) {$M$};
\node[box, right=of M] (Y) {$Y$};
\node[box, above=of M] (U) {$U$};
\draw[arr] (A) -- (M);
\draw[arr] (M) -- (Y);
\draw[arr] (U) -- (A);
\draw[arr] (U) -- (Y);
\end{tikzpicture}
\caption{Front-door graph: $A$ and $Y$ are confounded by $U$, but the effect of $A$ on $Y$ is fully mediated by observed $M$.}
\label{fig:ch9-frontdoor}
\end{figure}

\subsection{Mediator-Based Identification}

The graph says that the causal effect decomposes into two parts:
\[
  A\rightarrow M,\qquad M\rightarrow Y.
\]
The effect of $A$ on $M$ is unconfounded. The effect of $M$ on $Y$ can be
identified by adjusting for $A$, because the remaining noncausal path from $M$
to $Y$ goes through $A$.

\subsection{SWIGs for Front-Door Models}

The front-door argument can be expressed with counterfactual independencies:
\[
  M(a)\ci A,
\]
\[
  Y(m)\ci M\given A,
\]
\[
  Y(m)\ci M(a).
\]
The exclusion-type condition is
\[
  Y(a,m)=Y(m),
\]
which states that once the mediator is set to $m$, the value assigned to $A$
does not further affect $Y$.

\begin{figure}[H]
\centering
\begin{tikzpicture}[
  scale=.92,
  transform shape,
  box/.style={
    draw,
    rounded corners=2pt,
    align=center,
    minimum width=.9cm,
    minimum height=.65cm
  },
  fixed/.style={
    draw,
    rounded corners=2pt,
    fill=gray!12,
    align=center,
    minimum width=.9cm,
    minimum height=.65cm
  },
  arr/.style={-{Latex[length=2.1mm]}, thick}
]

% ---------- Row 1 ----------
\node[box]   (A1) at (0,0) {$A$};
\node[fixed] (a1) at (1.6,0) {$a$};
\node[box]   (M1) at (3.2,0) {$M(a)$};
\node[box]   (Y1) at (5.0,0) {$Y(a)$};
\node[box]   (U1) at (2.7,1.35) {$U$};

\draw[arr] (a1) -- (M1);
\draw[arr] (M1) -- (Y1);
\draw[arr] (U1) -- (A1);
\draw[arr] (U1) -- (Y1);

\node[right=1.0cm of Y1] {$M(a)\ci A$};

% ---------- Row 2 ----------
\node[box]   (A2) at (0,-2.6) {$A$};
\node[box]   (M2) at (1.6,-2.6) {$M$};
\node[fixed] (m2) at (3.2,-2.6) {$m$};
\node[box]   (Y2) at (5.0,-2.6) {$Y(m)$};
\node[box]   (U2) at (2.7,-1.25) {$U$};

\draw[arr] (A2) -- (M2);
\draw[arr] (m2) -- (Y2);
\draw[arr] (U2) -- (A2);
\draw[arr] (U2) -- (Y2);

\node[right=1.0cm of Y2] {$Y(m)\ci M\given A$};

% ---------- Row 3 ----------
\node[fixed] (a3) at (0,-5.2) {$a$};
\node[box]   (M3) at (1.6,-5.2) {$M(a)$};
\node[fixed] (m3) at (3.2,-5.2) {$m$};
\node[box]   (Y3) at (5.0,-5.2) {$Y(m)$};
\node[box]   (U3) at (2.7,-3.85) {$U$};

\draw[arr] (a3) -- (M3);
\draw[arr] (m3) -- (Y3);
\draw[arr] (U3) -- (Y3);

\node[right=1.0cm of Y3] {$Y(m)\ci M(a)$};

\end{tikzpicture}
\caption{SWIG-style views of the counterfactual independencies used in front-door identification.}
\label{fig:ch9-frontdoor-swig}
\end{figure}

\subsection{Identification Steps}

Start from the distribution of $Y(a)$:
\[
  p\{Y(a)\}=\sum_m p\{Y(a),M(a)=m\}.
\]
By the mediation condition, when $M(a)=m$,
\[
  Y(a)=Y(a,m)=Y(m).
\]
Thus
\[
  p\{Y(a)\}=\sum_m p\{Y(m),M(a)=m\}.
\]
Using the counterfactual independence $Y(m)\ci M(a)$,
\[
  p\{Y(a)\}=\sum_m p\{Y(m)\}p\{M(a)=m\}.
\]
The mediator distribution is identified because $M(a)\ci A$ and consistency:
\[
  p\{M(a)=m\}=p(M=m\given A=a).
\]
The remaining term is
\[
  p\{Y(m)\}=\sum_{a'}p\{Y(m)\given A=a'\}p(A=a').
\]
Using $Y(m)\ci M\given A$ and consistency,
\[
  p\{Y(m)\given A=a'\}=p(Y\given M=m,A=a').
\]

\subsection{Front-Door Adjustment Formula}

\begin{identificationbox}
For discrete $A$ and $M$, the front-door formula is
\[
  \E\{\Ya{a}\}
  =
  \sum_m
  \left[
    \sum_{a'} \E(Y\given M=m,A=a')\Prob(A=a')
  \right]
  \Prob(M=m\given A=a).
\]
\end{identificationbox}

The assigned value $a$ appears only in the mediator distribution
$\Prob(M=m\given A=a)$. The inner sum identifies the effect of setting the
mediator to $m$ on the outcome by averaging over the treatment distribution.

\subsection{Key Assumptions}

\begin{assumptionbox}
Front-door identification requires:
\begin{itemize}
  \item all directed effects of $A$ on $Y$ pass through $M$;
  \item no unmeasured confounding of $A$ and $M$;
  \item no unmeasured confounding of $M$ and $Y$ after conditioning on $A$;
  \item positivity for the relevant treatment and mediator levels.
\end{itemize}
\end{assumptionbox}

If there is a direct arrow $A\rightarrow Y$, then $Y(a,m)=Y(m)$ fails. If there
is unmeasured confounding between $A$ and $M$, then $p(M(a))$ is not identified
from $p(M\given A=a)$. If there is unmeasured confounding between $M$ and $Y$
within levels of $A$, then the mediator effect is not identified.

\subsection{Practical Implications}

The front-door formula is rarely justified automatically. It requires a mediator
that captures the entire causal effect of treatment and is itself sufficiently
well measured. Nevertheless, it is important because it proves that adjustment
is not the only route to nonparametric identification.

\subsection{Extensions and Software}

Modern identification algorithms generalize front-door reasoning to larger
graphs with hidden variables. Software can take a graph with observed and hidden
variables and return an identifying functional when one exists. Estimation then
requires fitting the conditional distributions appearing in that functional,
often with one-step or TMLE-style estimators.

\begin{figure}[H]
\centering
\begin{tikzpicture}[
  node distance=1.0cm,
  box/.style={draw, rounded corners=2pt, minimum width=.9cm, minimum height=.65cm},
  arr/.style={-{Latex[length=2.1mm]}, thick},
  bi/.style={<->, >=Latex, thick, dashed}
]
\node[box] (X) {$X$};
\node[box, right=of X] (A) {$A$};
\node[box, right=of A] (M) {$M$};
\node[box, right=of M] (L) {$L$};
\node[box, right=of L] (Y) {$Y$};
\draw[arr] (X) -- (A);
\draw[arr] (X) to[bend left=22] (M);
\draw[arr] (X) to[bend left=30] (Y);
\draw[arr] (A) -- (M);
\draw[arr] (M) -- (L);
\draw[arr] (L) -- (Y);
\draw[bi] (A) to[bend right=22] (Y);
\end{tikzpicture}
\caption{Extended front-door-type graphs can be identifiable even when no ordinary adjustment set exists.}
\label{fig:ch9-extended-frontdoor}
\end{figure}

\begin{implementationbox}
Graph-based software for hidden-variable causal models typically asks the user
to provide the observed variables, directed edges, hidden-variable or bidirected
edges, the treatment, and the outcome. The software then checks whether the
effect is identifiable and, when possible, returns an identifying functional and
estimators such as one-step or TMLE estimators. The graph specification is part
of the statistical analysis; changing the graph changes the estimand that can be
identified.
\end{implementationbox}

\begin{keyidea}
Front-door identification works by replacing an unidentifiable direct treatment
effect with two identifiable pieces: the effect of treatment on mediator and the
effect of mediator on outcome.
\end{keyidea}
